El sistema DSM convierte un conjunto de procesos independientes sobre Linux en un espacio de memoria compartida distribuida. Cada proceso que actúa como nodo aporta un bloque de memoria física, proyecta la misma vista virtual de 4 GiB en una dirección base común y resuelve los accesos a páginas remotas mediante fallos de página y el protocolo de coherencia MSI. La aplicación usa punteros normales: cuando toca una página que no tiene, el manejador de \texttt{SIGSEGV} registra el fallo, un despachador lo resuelve por TCP con el nodo custodio y el hilo reintenta la instrucción sin intervención del programa.

El sistema garantiza coherencia MSI por página con estados transitorios, un custodio determinista calculado con un anillo de dispersión consistente ponderado por memoria, tolerancia a fallos mediante franjas con paridad XOR, deltas y reconstrucción, y ejecución distribuida de kernels sobre particiones descritas en la vista. La membresía se mantiene con anuncios multicast y broadcast de respaldo, ranuras, épocas y latidos; la configuración compartida se resume en un hash que rechaza nodos incompatibles, mientras la memoria aportada es local a cada equipo. La herramienta \texttt{dsmctl} y el reporte HTML exponen la membresía, el directorio y el estado de las páginas.

La validación cubre la compilación sin avisos en cuatro perfiles, pruebas de extremo a extremo de coherencia, ejecución distribuida y reconstrucción con paridad, y un clúster de cuatro nodos que verifica la multiplicación de matrices contra una referencia secuencial.
